% 核心结论：两个时刻的空间合磁场
% 图原型：纯几何／对比。物理网络按本题节点、端口和符号构建。
\documentclass[tikz,border=7pt]{standalone}
\usepackage{fontspec,xeCJK,amsmath,amssymb}
\IfFontExistsTF{Times New Roman}{\setmainfont{Times New Roman}}{\setmainfont{TeX Gyre Termes}}
\IfFontExistsTF{Heiti SC}{\setCJKmainfont[AutoFakeBold=1.5]{Heiti SC}}{\setCJKmainfont[AutoFakeBold=1.5]{FandolHei-Regular}}
\usetikzlibrary{arrows.meta,calc,angles,quotes}
\definecolor{ink}{HTML}{243745}
\definecolor{blue}{HTML}{275DAB}
\definecolor{green}{HTML}{177C63}
\definecolor{red}{HTML}{BD4935}
\definecolor{purple}{HTML}{715393}
\begin{document}
\begin{tikzpicture}[font=\fontsize{9.5}{12}\selectfont,text=ink,
 v/.style={line width=1.1pt,-{Stealth[length=2.4mm,width=1.5mm]}},
 axis/.style={ink!65,line width=.45pt,-{Stealth[length=1.7mm,width=1.1mm]}},
 aux/.style={ink!40,densely dashed,line width=.5pt},
 lab/.style={fill=white,inner sep=2pt},
 card/.style={draw=blue!25,fill=blue!3,rounded corners=3pt,text width=6cm,align=left,inner sep=10pt}]
\path[use as bounding box] (-8,-3.35) rectangle (8,3.65);
\node[font=\fontsize{10.5}{13}\selectfont\bfseries,text width=15cm,align=center] at (0.00000,3.22000) {M6-076-P1　两个时刻的空间合磁场};
\begin{scope}[shift={(-4,0)}]
\draw[axis] (-2.50000,0.00000) -- (2.60000,0.00000);
\draw[axis] (0.00000,-2.05000) -- (0.00000,2.25000);
\node[] at (2.58000,-0.26000) {$x$};
\node[] at (-0.33000,2.20000) {$y$};
\node[] at (-0.18000,-0.20000) {$O$};
\draw[v,red] (0.00000,0.00000) -- (2.25000,0.00000);
\node[text=red] at (2.59450,0.36463) {$\mathbf B$};
\draw[aux] (0.00000,0.00000) -- (2.00000,0.00000);
\draw[v,blue] (0.00000,0.00000) -- (1.50000,0.00000);
\node[lab,text=blue] at (1.65000,-0.30000) {$\mathbf B_A$};
\draw[aux] (0.00000,0.00000) -- (-1.00000,1.73205);
\draw[v,green] (0.00000,0.00000) -- (0.37500,-0.64952);
\node[lab,text=green] at (0.90000,-1.00000) {$\mathbf B_B$};
\draw[aux] (0.00000,0.00000) -- (-1.00000,-1.73205);
\draw[v,purple] (0.00000,0.00000) -- (0.37500,0.64952);
\node[lab,text=purple] at (0.80000,1.00000) {$\mathbf B_C$};
\node[] at (0.00000,-2.50000) {$t=0$};
\end{scope}
\begin{scope}[shift={(4,0)}]
\draw[axis] (-2.50000,0.00000) -- (2.60000,0.00000);
\draw[axis] (0.00000,-2.05000) -- (0.00000,2.25000);
\node[] at (2.58000,-0.26000) {$x$};
\node[] at (-0.33000,2.20000) {$y$};
\node[] at (-0.18000,-0.20000) {$O$};
\draw[v,red] (0.00000,0.00000) -- (1.12500,1.94856);
\node[text=red] at (0.98147,2.42922) {$\mathbf B$};
\draw[aux] (0.00000,0.00000) -- (2.00000,0.00000);
\draw[v,blue] (0.00000,0.00000) -- (0.75000,0.00000);
\node[lab,text=blue] at (1.05000,-0.30000) {$\mathbf B_A$};
\draw[aux] (0.00000,0.00000) -- (-1.00000,1.73205);
\draw[v,green] (0.00000,0.00000) -- (-0.37500,0.64952);
\node[lab,text=green] at (-1.10000,0.80000) {$\mathbf B_B$};
\draw[aux] (0.00000,0.00000) -- (-1.00000,-1.73205);
\draw[v,purple] (0.00000,0.00000) -- (0.75000,1.29904);
\node[lab,text=purple] at (1.60000,1.10000) {$\mathbf B_C$};
\node[] at (0.00000,-2.50000) {$t=T/6$};
\end{scope}
\node[text width=15cm,align=center] at (0.00000,-3.00000) {三条彩色细矢量为各相瞬时磁场，红色为合场；共用长度尺度，合场大小恒为3Bm/2。};

\end{tikzpicture}\end{document}